\documentclass[10pt]{article}
\usepackage{icmlko}

\icmlmeta{IP 파운데이션 월드모델}{969}{버려진 것을 모른다고 불렀다}{앞 세 노트가 「통제」라 부른 열은 안 잰 것이 아니라 판이 결과를 보고 일부러 버린 것이었다 --- 되살려 보니 명제는 안 바뀐다}

\begin{document}

\twocolumn[
\icmlheader

\begin{abstract}
앞 노트(968)는 \texttt{wiki\_level} 이 \textbf{12 도메인 중 7 에서 마스크가 한 행도
안 선다}는 것을 잡고 그것을 \textbf{「모른다 --- 안 쟀다」}로 라벨했다. 그 앞 노트(966)는
같은 자리를 두고 \textbf{「애니에는 수준이 없다」}고 적었다. \textbf{둘 다 틀렸다.}
그 자리는 \emph{안 잰} 것이 아니라 \textbf{판의 설정 상수 한 줄이 일부러 버린} 것이다 ---
\texttt{lab/loop.py} 의 \texttt{WIKI\_DROP} 이 다섯 도메인에서 위키 축을 통째로 빼고,
\texttt{TREND\_DROP} 이 팝업에서 검색 축을 뺀다. 그 상수를 \textbf{메모리 안에서만} 비우고
다시 세니 애니의 수준은 \textbf{403 행 · 372 가지}가 살아난다(966 이 「없다」고 적은
바로 그 열이다). 판 전량으로는 \textbf{436 칸 중 18 칸 · 실측 2{,}052 행}이 그렇게
버려져 있고, \textbf{잰 행이 0 인 301 칸 중 나머지 283 만 진짜로 없다.}
\textbf{그래서 조항 59 에 셋째 갈래를 넣는다}: 「0 행」·「결측」·\textbf{「쟀는데 설정이
버렸다」}. 셋째를 첫째로 읽으면 인식론이 뒤집힐 뿐 아니라 \textbf{선택 편향이 안 보인다} ---
사료를 캐 보니 그 상수는 \textbf{유보 정답으로 채점한 판 $\rho$ 를 보고} 정해졌다
(노트 350 · $0.4859 \to 0.4931$ · 12 씨앗 중 10 승). \textbf{판에서는 옳은 결정이었다}
--- 사전등록도 12씨앗 대조도 원장 항목도 다 있다. \textbf{그러나 그 설정이 살아 있는 채로
명제의 통제 집합을 짜면 「어느 도메인에 통제가 있나」가 결과와 상관된 양으로 선택된다.}
\textbf{그러면 명제는 무너지는가 --- 아니다.} 그 상수를 비우고 규격 D 를
\textbf{50{,}000 뽑기}로 다시 재니 Holm$(\alpha{=}.05, m{=}10)$ 통과 수가
\textbf{두 이름 모두 불변}이다(긴 띠 $2\to2$ · 들뜸 $3\to3$). 개별 $\rho$ 는 움직인다
(팝업 $0.0676$ · 도서 $0.0496$) --- \textbf{문턱 근처가 아니어서 결론이 안 바뀐 것이지
되살린 통제가 무력해서가 아니다.} 곁들여 \textbf{게임을 닫았다}: \textbf{200{,}000 뽑기}에서
들뜸 이름 $p = 0.0055$ 로 문턱에서 \textbf{4.44 SE} 통과, 긴 띠 이름 $p = 0.0179$ 로
\textbf{39.33 SE} 탈락 --- \textbf{「2 냐 3 이냐」를 정한 것은 씨앗이 아니라 이름이다.}
\end{abstract}
]

\section{무엇을 정정하는가}

\claimx{966 의 「애니에는 수준이 없다」와 968 의 「모른다 --- 안 쟀다」는 둘 다 틀렸다 --- 애니의 수준은 403 행 · 372 가지가 있고, 판의 설정 상수가 그것을 버리고 있었다}

이 저장소의 판은 도메인마다 $(A, M, y, t)$ 를 들고 있고 $M[i,j]$ 는 「$i$ 행에서 $j$ 축을
쟀나」다. \texttt{lab/loop.py} 의 \texttt{WIKI\_DROP} 은 다섯 도메인(웹툰 · 팝업 · 애니 ·
시장팝업 · 도서)을 위키 축 dict 에서 \textbf{통째로 뺀다.} 그러면 \texttt{lab/harness.py}
가 \textbf{경고도 없이} 값 $0.5$ · 마스크 $0$ 으로 채운다 --- 경고는 그 도메인이 dict 에
\emph{있을} 때만 나는데 여기서는 아예 없기 때문이다.

그 상수를 \textbf{메모리 안에서만} 비우고(\texttt{L.WIKI\_DROP = ()}) 다시 세면:

\begin{center}\small
\begin{tabular}{lrrr}
도메인 & 현행 잰 행 & \textbf{비움 잰 행} & \textbf{값 가짓수} \\
\hline
애니 & 0 & \textbf{403} & \textbf{372} \\
시장팝업 & 0 & 83 & 81 \\
웹툰 & 0 & 52 & 52 \\
팝업 & 0 & 47 & 47 \\
도서 & 0 & 25 & 23 \\
\hline
영화 & 0 & \textbf{0} & 0 \\
펀딩 & 0 & \textbf{0} & 0 \\
\end{tabular}
\end{center}

\textbf{영화와 펀딩만 진짜로 없다.} 영화는 위키 개체 목록에 아예 없고, 펀딩은
\texttt{lab/wikiaxes.py} 의 \texttt{if np.isfinite(raw).sum() < 20: continue} 게이트에
걸린다 --- \textbf{그것은 「버린 것」이 아니라 「진짜로 못 잰 것」이다.}

판 전량 $(12 \times \text{열})$ \textbf{436 칸}을 세 갈래로 세면
\textbf{산다 135 · 쟀는데 설정이 버렸다 18 · 0 행 283} 이다. 버린 18 칸은
위키 5 도메인 $\times$ 3 축 $+$ 검색 1 도메인(팝업) $\times$ 3 축이고, 버린 실측 행은
\textbf{2{,}052} 다.

\claimx{그 설정 상수는 유보 정답으로 채점한 판 $\rho$ 를 보고 정해졌다 --- 판에서는 옳은 결정이고 명제 측정에서는 옳지 않다}

\texttt{WIKI\_DROP} 이 왜 있는지 \textbf{지우기 전에} 캤다. 이력이 GitHub 이관 스냅샷
하나로 뭉개져 있어서 보존 가지에서 다시 셌더니 도입 커밋이 하나 나온다 --- \textbf{노트 350}
(2026-08-01). 같은 커밋에 \textbf{사전등록 파일}과 \textbf{원장 항목}이 함께 들어 있고,
근거는 한 줄로 요약된다: \textbf{판 $\rho$ 가 $0.4859 \to 0.4931$ 로 올랐다}
($\Delta = +0.0033$ · 12 씨앗 중 \textbf{10} 에서 이겼다). 누출 방지도 표본 부족도
아니다 --- \textbf{절제(ablation)에 의한 축 선택}이다.

\textbf{이 저장소의 관례로는 모범적인 결정이다.} 사전등록에 「미리 적는 실패」까지 적혀 있다.
\textbf{그러나 판정용 자로 물려쓰면 두 가지가 어긋난다.} 첫째, \textbf{물음이 다르다} ---
판은 「이 열이 예측을 돕나」를 묻고 버렸는데, 통제는 「이 열이 교란인가」를 묻는다.
\textbf{예측을 안 돕는 열도 교란일 수 있다.} 둘째이자 더 무거운 것은,
\textbf{판 $\rho$ 가 유보 정답으로 채점한 수}라는 것이다. 그러면 「어느 도메인에 통제가
살아 있나」가 \textbf{결과와 상관된 양으로 선택된} 것이고, 966$\sim$968 이 막으려던 바로
그 종류의 오염이 \textbf{통제 집합의 구성}에 들어와 있다.

\claimx{그 상수를 비우고 50,000 뽑기로 다시 재도 Holm 통과 수는 두 이름 모두 불변이다 --- 명제는 이 정정에서 살아남는다}

규격 D(동률평균 순위 · 자기 열의 앞항 통제 · 두 띠 모두 실자료)를
\textbf{Holm$(\alpha{=}.05, m{=}10)$ · 50{,}000 뽑기 · 씨앗 12345} 로 두 팔에 걸었다.

\begin{center}\small
\begin{tabular}{lrr}
 & 현행 & \textbf{\texttt{WIKI\_DROP}$=()$} \\
\hline
긴 띠 이름 통과 & 2 & \textbf{2} \\
들뜸 이름 통과 & 3 & \textbf{3} \\
\end{tabular}
\end{center}

통제가 되살아난 다섯 도메인에서 $|\Delta\rho|$ 는 팝업 $0.0676$ · 도서 $0.0496$ ·
시장팝업 $0.0138$ · \textbf{애니 $0.0072$} · 웹툰 $0.0003$ 이다. \textbf{크게 움직이는
자리가 있는데도 결론이 안 바뀐 이유는 그 자리들이 문턱 근처가 아니기 때문이다} ---
되살린 통제가 무력해서가 아니다. \textbf{이것은 재 봐야 아는 것이었다.}

\claimx{「Holm 통과 2 냐 3 이냐」를 정한 것은 씨앗이 아니라 이름이다 --- 200,000 뽑기에서 두 이름이 39 SE 로 갈린다}

968 은 게임의 $p$ 가 씨앗에 따라 $0.0060 \leftrightarrow 0.0090$ 을 오간다고 신고했고
그것을 \textbf{「결정이 잡음 안에 산다」}로 읽었다. 뽑기를 \textbf{200{,}000} 으로 올리면
잡음이 걷힌다:

\begin{center}\small
\begin{tabular}{lrrr}
이름 & $\rho$ & $p$ & 문턱에서 \\
\hline
긴 띠 & $-0.1976$ & $0.01791$ & \textbf{39.33 SE 탈락} \\
들뜸 & $+0.2311$ & $0.00552$ & \textbf{4.44 SE 통과} \\
\end{tabular}
\end{center}

\textbf{둘 다 굳었다.} 씨앗이 답을 흔든 것이 아니라 \textbf{2{,}000 뽑기가 모자랐던} 것이고,
답을 정하는 것은 \textbf{어느 이름으로 재나} 다. \texttt{excite} $=$ \texttt{recent}
$-$ \texttt{long} 이라 두 이름은 같은 양의 두 얼굴인데, 순위를 도메인 전체에서 매기고
잔차화를 그 뒤에 하므로 \textbf{선형 항등식이 순위를 지나며 깨진다.}
\textbf{사전등록이 정한 이름은 긴 띠이고, 그 이름으로는 2 다.}

\section{남은 것}

\textbf{이 노트가 못 한 것을 적는다.} 되살린 통제로 \textbf{판}(예측 성능)을 다시 돌리지
않았다 --- 축 목록이 전 도메인 이름의 교집합이라 도메인마다 다른 열을 빼면 공통 축이
무너지는 문제(968 이 잡았다)를 아직 안 풀었다. 그리고 여기 실린 것은 전부
\textbf{유보 위의 연관}이지 \textbf{예측 성능}이 아니다 --- 「긴 띠가 결과 정보를 담는다」와
「긴 기억이 도움이 된다」는 둘이고, 살아남은 기전의 부호는 \textbf{뒤쪽을 지지하지 않는다.}

\end{document}
